% An introductory chapter contextualising the project and the specific
% module or set of tasks worked on.


\section{Project Objectives and Requirements}

% Describe the project objectives and requirements. Include a brief outline
% of the methods and processes used.

The main objective of the project is to create an application for heat production optimization for a district heating utility,
optimizing heat production for the lowest cost, 
while allowing for two scenarios with different production assets, winter and summmer periods and enabling for maintenance periods.
The system must be separated into five modules: asset manager, source data manager, result data manager, optimizer and data visualization. 

The project has been divded into tasks for each sprint and distributed among all group memebers,
while having frequent ceremonies to ensure everyone keeps up with the workflow.


\section{Problem Statement}

% State the problem your project addresses. This should be a clear,
% concise formulation of the issue you set out to solve.

District heating utilities need to provide heat to all buildings in the districts heating network.
Providing heat to all buildings can become expensive and unsustainable, especially during winter periods where the heat demand is high.
An application can help manage the heating utilities by optimizing heat production for the lowest cost
and allowing for maintenance while simultaneously displaying the optimization results.